Con base en la Tab. \ref{tab:medidas} es posible afirmar que los amperimetros
empleados en los dos circuitos reportaron un valor altamente exacto, con una
desviación del órden de \qty{0,01}{\milli\ampere}. Sin embargo, el amperimetro
del circuito \ref{fig:amp-25} obtuvo el valor m\'as cercano de los dos
debido a que la resistencia interna del instrumento es mucho menor que la
resistencia $R$ total del circuito, en comparación con el circuito
\ref{fig:amp-125}. En consecuencia, el amperímetro del circuito \ref{fig:amp-25}
drena menos voltaje y modifica menos el circuito de estudio.

Por otro lado, los circuitos medidos con el voltímetro presentaron, en general,
un mayor error relativo; específicamente, el circuito \ref{fig:volt-100k-12}
con las resistencias $R_1$ y $R_2$ del órden de \qty{100}{\kilo\ohm} tuvo la
menor exactitud en el intervalo $V^{max}=12$. En contraste con el
amperímetro, el volt\'imetro tiene mayor exactitud cuando su
resistencia interna es mucho mayor que la resistencia del circuito que se esta
analizando, debido a que esta configuración drena poca corriente.

Además, aquellos circuitos analizados con el intervalo de voltaje
$V^{max}=4$ presentaron medidas alejadas del valor real, lo cual es un
resultado esperado. Dado que al medir el circuito con un intervalo de voltaje menor
implica una menor resistencia interna del instrumento, este drena más corriente,
afectando el resultado.

Por último, realizar los circuitos en una simulación elimina múltiples problemas, p.ej,
no se consideran la incertidumbre de medida, ni las asociadas a las resistencias y
fuentes de voltaje. Similarmente, tampoco se considera la resistencia aportada por
cables ni la interna de la fuente. Por esta razón, los resultados fueron los esperados
por la teoria sin ser afectados por factores desconocidos.
